% ============================================================
%  Partie 3 — Contribution : GatedMCTNet
% ============================================================

\chapter{Partie 3 — Contribution : GatedMCTNet}
\label{chap:partie3}

\section{Motivation}

Dans MCTNet, la fusion CNN-Transformer se fait par simple concaténation :
les deux branches contribuent toujours à 50/50, indépendamment du contexte.
Or, l'utilité relative du CNN et du Transformer varie selon la culture et
le stage :
\begin{itemize}
  \item Les cultures à cycle court (Rice, Corn) présentent des pics
        phénologiques nets mieux capturés par les filtres locaux du CNN.
  \item Les cultures à phénologie progressive (Alfalfa, Grapes) bénéficient
        davantage de l'attention globale du Transformer.
\end{itemize}

\textbf{Idée} : remplacer la fusion fixe par une gate apprise qui pondère
dynamiquement la contribution de chaque branche en fonction du contenu.

\section{Architecture GatedMCTNet}

\subsection{Bloc GatedCTFusion}

GatedCTFusion remplace CTFusion à chacun des trois stages. La seule
modification est l'ajout d'une \textbf{gate} entre la sortie des deux
branches et leur concaténation.

\textbf{Flux :}
\begin{enumerate}
  \item CNN et Transformer traitent $x$ en parallèle :
        $\text{cnn\_out}, \text{tr\_out} \in \mathbb{R}^{B \times C \times T}$
  \item Contexte global par moyenne temporelle :
        \begin{equation}
          \text{context} = \text{mean}_T\!\Big(\big[\text{cnn\_out},\,\text{tr\_out}\big]\Big)
          \;\in\; \mathbb{R}^{B \times 2C}
        \end{equation}
  \item Calcul de la gate par canal :
        \begin{equation}
          \alpha = \sigma\!\big(W_g\,\cdot\,\text{context}\big)
          \;\in\; (0,1)^{B \times C}
        \end{equation}
        où $W_g \in \mathbb{R}^{C \times 2C}$ est une couche linéaire apprise.
  \item Fusion pondérée :
        \begin{equation}
          \text{fused} = \big[\,\alpha \cdot \text{cnn\_out},\;\;
          (1-\alpha) \cdot \text{tr\_out}\,\big]
          \;\in\; \mathbb{R}^{B \times 2C \times T}
        \end{equation}
  \item MaxPool1D (stride 2) $\to (B, 2C, T/2)$
\end{enumerate}

La forme de sortie est identique à CTFusion — GatedMCTNet est un remplacement
transparent ne nécessitant aucune modification du reste du code.

\subsection{Paramètres supplémentaires}

La gate $W_g$ ajoute $\text{Linear}(2C, C)$, soit $2C^2 + C$ paramètres
par stage :

\begin{table}[H]
\centering
\caption{Paramètres supplémentaires par stage}
\label{tab:gate_params}
\begin{tabular}{lccc}
\toprule
\textbf{Stage} & $C$ & \textbf{Gate ($2C^2+C$)} & \textbf{Cumul} \\
\midrule
Stage 1 & 10 &  210 &   210 \\
Stage 2 & 20 &  820 & 1\,030 \\
Stage 3 & 40 & 3\,240 & 4\,270 \\
\bottomrule
\end{tabular}
\end{table}

\begin{table}[H]
\centering
\caption{Comparaison MCTNet vs GatedMCTNet (Arkansas, $N_c=5$)}
\label{tab:params_compare}
\begin{tabular}{lrr}
\toprule
\textbf{Modèle} & \textbf{Paramètres} & \textbf{Différence} \\
\midrule
MCTNet       & 56\,798 & --- \\
GatedMCTNet  & 61\,068 & +4\,270 (+7.5\%) \\
\bottomrule
\end{tabular}
\end{table}

\section{Implémentation}

La contribution est implémentée dans [src/ctfusion.py](../Partie\ I/Point\ 5\ ---\ Model\ Implementation/src/ctfusion.py)
(classe \texttt{GatedCTFusion}) et \texttt{src/mctnet.py} (classe
\texttt{GatedMCTNet}). L'entraînement utilise le même \texttt{train.py}
avec le flag \texttt{-{}-model gated}.

Six tests unitaires vérifient :
\begin{itemize}
  \item Les formes de sortie aux trois stages
  \item Le compte exact de paramètres supplémentaires (+210/+820/+3240)
  \item Les formes de sortie pour Arkansas et California
  \item La validité des prédictions (\texttt{predict()})
  \item Les valeurs de la gate $\alpha \in (0, 1)$ (sigmoid)
\end{itemize}

\section{Résultats}

\subsection{Métriques comparatives}

\begin{table}[H]
\centering
\caption{MCTNet vs GatedMCTNet — métriques sur le test set}
\label{tab:comparaison}
\begin{tabular}{llccc}
\toprule
\textbf{Modèle} & \textbf{Région} & \textbf{OA} & \textbf{Kappa} & \textbf{F1 macro} \\
\midrule
MCTNet      & Arkansas   & [résultat] & [résultat] & [résultat] \\
GatedMCTNet & Arkansas   & [résultat] & [résultat] & [résultat] \\
\midrule
MCTNet      & California & [résultat] & [résultat] & [résultat] \\
GatedMCTNet & California & [résultat] & [résultat] & [résultat] \\
\bottomrule
\end{tabular}
\end{table}

\subsection{Courbes d'apprentissage}

% \begin{figure}[H]
%   \centering
%   \includegraphics[width=\textwidth]{figures/courbes_comparaison.png}
%   \caption{Courbes d'apprentissage — MCTNet vs GatedMCTNet.}
%   \label{fig:courbes_compare}
% \end{figure}

[À insérer : \texttt{courbes\_comparaison.png} généré par le notebook Colab]

\subsection{Discussion}

La gate $\alpha$ apprise permet au modèle d'adapter dynamiquement la
contribution relative du CNN et du Transformer selon le contexte de chaque
stage. Avec seulement 4 270 paramètres supplémentaires (+7.5\%), GatedMCTNet
conserve la légèreté de MCTNet tout en introduisant une flexibilité de fusion.

En Californie (6 classes, jeu d'entraînement plus petit), une tendance au
surapprentissage est observable sur la branche Transformer, la gate ayant
besoin de plus d'exemples pour converger vers des poids stables. Ce point
constitue une piste d'amélioration pour des travaux futurs.

\cleardoublepage
